\documentclass[12pt,a4paper]{article}

\usepackage[margin=1in]{geometry}
\usepackage[T1]{fontenc}
\usepackage[utf8]{inputenc}
\usepackage{lmodern}
\usepackage{enumitem}
\usepackage{hyperref}

\title{Documentation (20\%)\\Jobra Hospital Management System}
\author{}
\date{\today}

\begin{document}
\maketitle

\section{Code documentation}
\begin{itemize}[leftmargin=*, itemsep=0.4em]
  \item \textbf{Readable structure over heavy comments}: The project relies mainly on clear file/folder organization and descriptive naming (controllers, routes, services) rather than extensive inline comments.
  \item \textbf{Inline intent where helpful}: Some files include brief intent comments (e.g., route grouping, ``test route'', or local query helpers), which helps onboarding.
  \item \textbf{Improvement note}: Adding brief docstrings/comments for key controllers, expected request bodies, and role requirements would make the codebase easier to maintain for new contributors.
\end{itemize}

\section{Technical documentation}
\begin{itemize}[leftmargin=*, itemsep=0.4em]
  \item \textbf{Comprehensive README}: The root \texttt{README.md} documents features, tech stack, high-level structure, database overview, setup instructions for backend/frontend, and suggested future improvements.
  \item \textbf{Database schema included}: The \texttt{database.sql} file provides the full relational schema for the \texttt{jobra\_hospital} database, serving as a concrete technical reference.
  \item \textbf{Configuration guidance}: The README explains required environment variables for database connectivity and highlights security considerations around JWT secrets.
\end{itemize}

\section{User documentation}
\begin{itemize}[leftmargin=*, itemsep=0.4em]
  \item \textbf{Role-based user flows}: The README outlines patient/doctor/admin flows (login, booking, report writing, complaint handling), which works as lightweight user documentation.
  \item \textbf{UI discoverability}: The dashboard-based navigation (admin/doctor/patient) reduces the need for long user manuals by keeping actions grouped per role.
  \item \textbf{Improvement note}: A short ``How to use'' guide with screenshots for each role (admin portal, doctor portal, patient portal) would strengthen user-facing documentation.
\end{itemize}

\section{API documentation}
\begin{itemize}[leftmargin=*, itemsep=0.4em]
  \item \textbf{Documented API surface (implicit)}: API prefixes are consistent under \texttt{/api/...} and the codebase groups endpoints by domain (auth/profile/admin/patient/doctor/appointment/report/complaint).
  \item \textbf{Auth contract}: The frontend clearly sends the JWT in the \texttt{Authorization} header via an Axios interceptor, and the backend middleware verifies it to populate \texttt{req.user}.
  \item \textbf{Improvement note}: A formal API reference (OpenAPI/Swagger or a markdown endpoint table) listing endpoints, required role, request/response schemas, and error codes would make integration and testing easier.
\end{itemize}

\end{document}

